\subsection{Non-injective projection versus superdeterminism and free choice}
  \label{subsec:free-choice}

  A frequent objection to any violation of statistical independence in Bell-type analyses
  is that it implicitly relies on superdeterministic assumptions, namely correlations
  between measurement settings and underlying physical states.
  In such scenarios, the apparent violation of Bell inequalities is attributed to a
  conspiratorial coordination between experimental choices and microscopic conditions.

  The mechanism discussed here is fundamentally different.
  No correlation is postulated between measurement settings and underlying
  configurations, and no restriction is imposed on the freedom of experimental choices
  at the level of observable descriptions.
  Measurement settings are assumed to be freely and independently chosen within the
  effective description, in full accordance with standard experimental practice.

  The loss of statistical independence arises instead from the non-injective structure
  of the mapping from underlying configurations to observable outcomes.
  Such non-injectivity arises generically in regimes where the effective description
  has a finite capacity to encode distinctions present at the underlying level.
  Because distinct underlying configurations are identified at the effective level,
  conditioning on observable variables does not isolate independent subsystems.
  The resulting failure of probabilistic factorization is therefore structural and
  geometric in origin, rather than causal or conspiratorial.

  In this sense, the present framework preserves operational free choice while violating
  the assumptions required for Bell-type factorizations.
  The breakdown of statistical independence reflects intrinsic limitations of
  finite-capacity effective descriptions, not a fine-tuned coordination of initial conditions.
